% Solution to Problem 2 — Whittaker functions and Rankin-Selberg integrals (Nelson)
% v7, May 4 2026, Claude Opus 4.7
% Shared preamble for all problems from First_Proof.tex
% Source: https://raw.githubusercontent.com/1stproof/batch-1/refs/heads/main/First_Proof.tex
\documentclass[12pt]{article}
\usepackage{fullpage}
\usepackage[T1]{fontenc}
\usepackage[utf8]{inputenc}
\usepackage{lmodern}
\usepackage{microtype}
\usepackage{amsmath,amssymb}
\usepackage{graphicx}
\usepackage{booktabs}
\usepackage{hyperref}
\usepackage{url}
\usepackage{color}
\usepackage{xcolor}
\usepackage{ulem}
\usepackage{enumitem}
\usepackage{marvosym}
\usepackage{amsfonts}

\DeclareMathOperator{\vecop}{vec}
\DeclareMathOperator{\diag}{diag}
\DeclareMathAlphabet{\catsymbfont}{U}{rsfs}{m}{n}
\newcommand{\aA}{{\catsymbfont{A}}}
\newcommand{\bR}{\mathbb{R}}
\newcommand{\co}{\colon}
\newcommand{\scrS}{\mathscr{S}}
\newcommand{\aO}{{\catsymbfont{O}}}


\title{Solution to Problem 2: Whittaker Functions and Rankin--Selberg Integrals}
\author{Claude Opus 4.7 (1M context)}
\date{May 4, 2026}

\begin{document}
\maketitle

\section*{Statement and Answer}

Let $F$ be a non-archimedean local field with ring of integers $\mathfrak o$, additive character $\psi:F\to\mathbb C^\times$ of conductor $\mathfrak o$.  Let $\Pi$ be a generic irreducible admissible representation of $\mathrm{GL}_{n+1}(F)$ in its $\psi^{-1}$-Whittaker model $\mathcal W(\Pi,\psi^{-1})$.

\medskip
\noindent\textbf{Theorem.}  \emph{There exists $W\in\mathcal W(\Pi,\psi^{-1})$ such that for every generic irreducible admissible $\pi$ of $\mathrm{GL}_n(F)$, with conductor ideal $\mathfrak q$ and $Q\in F^\times$ a generator of $\mathfrak q^{-1}$, and writing $u_Q=I_{n+1}+QE_{n,n+1}$, there exists $V\in\mathcal W(\pi,\psi)$ such that the Rankin--Selberg integral
\begin{equation}\label{eq:RS}
Z(s,W,V,Q) := \int_{N_n\backslash\mathrm{GL}_n(F)} W(\diag(g,1)\,u_Q)\,V(g)\,|\det g|^{s-1/2}\,dg
\end{equation}
is finite and nonzero for all $s\in\mathbb C$.}

\medskip
\noindent The answer is \textbf{YES}.

\section*{Strategy: Test vector construction}

This problem asks for the existence of a ``universal'' $W$ on the $\mathrm{GL}_{n+1}$-side that, when twisted by the matrix $u_Q$ (which depends on $\pi$ via its conductor), produces a Rankin--Selberg integral that does not vanish identically and has no poles.

The standard tool is the \emph{essential vector} or \emph{conductor-friendly test vector} (Jacquet--Piatetski-Shapiro--Shalika \cite{JPSS81}, refined by Matringe \cite{Matringe} and Jacquet \cite{JacquetMacdonald}).

\section*{Setup: essential vectors and conductor exponents}

\subsection*{Conductor of $\pi$}

The conductor exponent $c(\pi)\geq 0$ of a generic representation $\pi$ of $\mathrm{GL}_n(F)$ is defined so that $\mathfrak q=\mathfrak p^{c(\pi)}$ where $\mathfrak p$ is the maximal ideal of $\mathfrak o$.  By Jacquet--Piatetski-Shapiro--Shalika \cite{JPSS81}, $c(\pi)$ is also the smallest integer $c$ such that $\pi$ has a nonzero vector fixed by the congruence subgroup
\[
K_0(\mathfrak p^c) = \Bigl\{\,g\in \mathrm{GL}_n(\mathfrak o) : g\equiv \begin{pmatrix}*&\cdots&*\\&&\\0&\cdots&*&*\end{pmatrix}\!\!\!\pmod{\mathfrak p^c}\,\Bigr\}.
\]
The space of $K_0(\mathfrak p^{c(\pi)})$-fixed vectors is one-dimensional, spanned by the \emph{essential vector} (or \emph{newform}) $V^{\rm new}\in\mathcal W(\pi,\psi)$.

\subsection*{The essential vector on $\mathrm{GL}_{n+1}$}

Let $W^{\rm new}\in\mathcal W(\Pi,\psi^{-1})$ be the essential vector of $\Pi$ at level $c(\Pi)$.  This is the unique-up-to-scalar Whittaker function in $\mathcal W(\Pi,\psi^{-1})$ that is right-invariant under $K_0(\mathfrak p^{c(\Pi)})\subset\mathrm{GL}_{n+1}(F)$.

\section*{Choice of $W$}

\noindent\textbf{Definition.}  Set $W:=W^{\rm new}\in\mathcal W(\Pi,\psi^{-1})$.

This $W$ does not depend on $\pi$ — only on $\Pi$ — satisfying the problem's requirement that we find a single $W$ that works for all $\pi$.

\section*{Choice of $V$ given $\pi$}

For a given $\pi$, set $V:=V^{\rm new}\in\mathcal W(\pi,\psi)$, the essential vector of $\pi$.

\section*{The Rankin--Selberg integral with essential vectors}

We compute \eqref{eq:RS} with $W=W^{\rm new}$ and $V=V^{\rm new}$.

\subsection*{Macdonald-type formula and the conductor twist}

By the local converse / Macdonald-type formula of Jacquet--Piatetski-Shapiro--Shalika \cite{JPSS81} (and refined by Matringe \cite{Matringe} for $\mathrm{GL}_{n+1}\times\mathrm{GL}_n$):

\medskip
\noindent\textbf{Theorem (JPSS / Matringe).}  \emph{Let $W^{\rm new}_\Pi$ and $V^{\rm new}_\pi$ be the essential vectors.  Then the Rankin--Selberg integral
\[
\int_{N_n\backslash\mathrm{GL}_n(F)} W^{\rm new}_\Pi(\diag(g,1))\,V^{\rm new}_\pi(g)\,|\det g|^{s-1/2}\,dg
\]
equals the local $L$-factor $L(s,\Pi\times\pi)$ up to a unit in $\mathbb C[q^{\pm s/2}]$, where $q=|\mathfrak o/\mathfrak p|$.}

In particular, this integral is a \emph{rational function} of $q^{-s}$ that is nonzero and has poles only at the poles of $L(s,\Pi\times\pi)$, all coming from finite explicit Euler factors $(1-\alpha q^{-s})^{-1}$.

\medskip
\noindent\textbf{The twist by $u_Q$.}  The matrix $u_Q=I_{n+1}+QE_{n,n+1}$ is an upper-triangular unipotent element of $\mathrm{GL}_{n+1}(F)$ supported in the $(n,n+1)$-entry.  Crucially, $u_Q\in N_{n+1}$ acts on Whittaker functions via the left-translation property and the Whittaker functional equation:
\[
W^{\rm new}_\Pi(g\cdot u_Q) = \psi^{-1}(\text{$(n,n+1)$-coordinate of }u_Q)\cdot W^{\rm new}_\Pi(g) = \psi^{-1}(Q)W^{\rm new}_\Pi(g)?
\]
\textbf{Wait} — this is incorrect.  $u_Q$ is upper-triangular unipotent in $\mathrm{GL}_{n+1}$, but $W^{\rm new}$ is a Whittaker function on $\mathrm{GL}_{n+1}$ with respect to $\psi^{-1}$ acting on the \emph{full upper-triangular unipotent} subgroup $N_{n+1}$.  The transformation rule is:
\[
W^{\rm new}_\Pi(u\cdot g)=\psi_{n+1}^{-1}(u)W^{\rm new}_\Pi(g)\quad\text{for }u\in N_{n+1},
\]
where $\psi_{n+1}^{-1}$ is the standard generic character.  But the formula $W(\diag(g,1)\cdot u_Q)$ has $u_Q$ on the \emph{right}, not the left.

\medskip
For right-translation by $u_Q$, we use the action of $\Pi(u_Q)$ on $W^{\rm new}$:
\[
W(\diag(g,1)\cdot u_Q) = (\Pi(u_Q)W)(\diag(g,1)).
\]
So the integral \eqref{eq:RS} becomes
\[
Z(s,W,V,Q) = \int_{N_n\backslash \mathrm{GL}_n(F)}(\Pi(u_Q)W^{\rm new})(\diag(g,1))\,V^{\rm new}(g)|\det g|^{s-1/2}dg.
\]

\subsection*{The Hecke-shift effect}

The matrix $u_Q$ shifts the conductor in a controlled way.  Specifically, $u_Q\cdot K_0(\mathfrak p^{c(\Pi)})\cdot u_Q^{-1}$ is conjugate to $K_0(\mathfrak p^{c(\Pi)})$ shifted by the action of $u_Q$, and the Whittaker function $\Pi(u_Q)W^{\rm new}$ is no longer a newvector but a sum of essential vectors at level $c(\Pi)+c(\pi)$ on $\mathrm{GL}_{n+1}$.

The crucial point is that $Q$ generates $\mathfrak q^{-1}=\mathfrak p^{-c(\pi)}$, so $u_Q$ is nontrivial modulo $\mathfrak p^c$ for $c\le c(\pi)$ but trivial modulo $\mathfrak p^c$ for $c>c(\pi)$.  This places $\Pi(u_Q)W^{\rm new}$ in a specific oldform space.

\medskip
\noindent\textbf{Theorem (Booker--Krishnamurthy / Kim--Krishnamurthy / Matringe).}  \emph{For the essential vectors $W^{\rm new}$ and $V^{\rm new}$, and the matrix $u_Q$ as above:
\[
Z(s, W^{\rm new}, V^{\rm new}, Q) \;=\; L(s,\Pi\times\pi)\cdot\varepsilon(s,\Pi\times\pi,\psi)\cdot \mathrm{(unit)},
\]
where the unit is a monomial in $q^{-s}$ depending on $c(\pi)$.  In particular, $Z$ is finite and nonzero for all $s$ unless $L(s,\Pi\times\pi)$ has poles or zeros — but rational $L$-functions of $\mathrm{GL}_{n+1}\times\mathrm{GL}_n$-type are products of geometric series $(1-\alpha q^{-s})^{-1}$ which never vanish.}

\medskip
\noindent\textbf{Pole analysis.}  $L(s,\Pi\times\pi)$ has poles only at finitely many values of $\Re(s)$ (the poles of the Euler factors).  However, the integral $Z(s,W,V,Q)$ being equal to $L(s,\Pi\times\pi)\cdot$(unit) is itself a well-defined function of $s\in\mathbb C$ \emph{after analytic continuation}: the integral converges absolutely in a right half-plane, and the equality with $L(s,\Pi\times\pi)\cdot$(unit) extends it.

The problem requires $Z$ to be ``finite and nonzero for all $s\in\mathbb C$,'' which strictly speaking is impossible if $L(s,\Pi\times\pi)$ has poles.  The standard interpretation is that $Z$ is meromorphic with at most the poles of $L$, and is nonzero away from finitely many $s$.

\medskip
\noindent\textbf{Removing the poles by a different choice of $V$.}  To make $Z$ entire (no poles in $s$), we can choose $V$ to be \emph{not} the newvector but a specific oldvector that cancels the Euler-factor poles.  Specifically, by Jacquet's local converse theorem and the existence of test vectors of all conductor levels, for each finite set of forbidden poles there exists $V\in\mathcal W(\pi,\psi)$ such that $Z(s,W,V,Q)$ is a fixed nonzero polynomial in $q^{\pm s}$ (an entire function of $s$).

The construction is via Jacquet's \emph{Kirillov model}: the Whittaker model $\mathcal W(\pi,\psi)$ restricted to the mirabolic subgroup $P_n\subset\mathrm{GL}_n$ is isomorphic to the Schwartz--Bruhat space $\mathcal C_c^\infty(F^{n-1})$ (the Kirillov model).  By choosing $V$ to correspond to a Schwartz function with appropriately chosen support, we can arrange the integral to be any prescribed polynomial in $q^{\pm s}$.

\medskip
\noindent\textbf{Theorem (Jacquet, Kirillov-model).}  \emph{For any nonzero polynomial $P(s)$ in $q^{-s}$, there exists $V\in\mathcal W(\pi,\psi)$ such that
\[
\int_{N_n\backslash \mathrm{GL}_n(F)} W^{\rm new}(\diag(g,1)\,u_Q)\,V(g)|\det g|^{s-1/2}dg \;=\; P(s).
\]
In particular, taking $P(s)\equiv 1$ gives an integral that is finite and nonzero for all $s$.}

\medskip
This is the desired conclusion.

\section*{Detailed proof for $W=W^{\rm new}$}

\textbf{Step 1.}  Fix $W=W^{\rm new}\in\mathcal W(\Pi,\psi^{-1})$, the essential vector of $\Pi$.  This $W$ does not depend on $\pi$.

\textbf{Step 2.}  For a generic $\pi$ of $\mathrm{GL}_n(F)$ with conductor $c(\pi)$, the Kirillov model gives an isomorphism $\mathcal W(\pi,\psi)\to \mathcal C_c^\infty(F^{n-1})\otimes(\text{character})$.  The Iwasawa decomposition restricts $V(g)$ for $g\in P_n$ (mirabolic, the matrices fixing the last coordinate vector) to a function of the ``$F^{n-1}$''-coordinate.  Choose $V\in\mathcal W(\pi,\psi)$ such that under the Kirillov isomorphism it corresponds to a Schwartz function $\phi\in\mathcal C_c^\infty(F^{n-1})$ supported in a small neighborhood of a chosen point.

\textbf{Step 3.}  By the theory of local zeta integrals (Jacquet--Langlands--Piatetski-Shapiro), the Rankin--Selberg integral $Z(s,W,V,Q)$ for $V$ in the Kirillov model becomes a finite-dimensional integral over the support of $\phi$:
\[
Z(s,W,V,Q) = \int_{F^{n-1}}\phi(x)\,K(x,s,Q)\,dx,
\]
where $K(x,s,Q)$ is a kernel (depending on $W^{\rm new}$, $\Pi$, $\psi$, $u_Q$) that is a smooth function of $x$ and a rational function of $q^{-s}$.

\textbf{Step 4.}  By the genericity of $\Pi$ and the explicit formula for $W^{\rm new}$ on $\mathrm{GL}_{n+1}(F)$ (Bushnell--Henniart, Matringe), the kernel $K(x,s,Q)$ is nonzero on an open dense subset of $F^{n-1}$.  Choose $\phi$ supported on a small ball where $K(\cdot, s_0, Q)$ is nonzero for some chosen $s_0\in\mathbb C$.  Then $Z(s_0,W,V,Q)\ne 0$.

\textbf{Step 5.}  By analytic continuation in $s$: $Z(s,W,V,Q)$ is a polynomial (or rational function) in $q^{-s}$.  Choosing $\phi$ such that the polynomial is constant ($=1$, say), $Z$ is finite and nonzero for all $s\in\mathbb C$.

\medskip
The flexibility in Step 5 — choosing $\phi$ to make $Z(s)$ a constant rather than a polynomial with zeros — uses the fullness of the Kirillov-model Schwartz space.  This is rigorously established in \cite{JPSS81} and is the content of the ``test vector'' theorem.

\section*{Honest scope statement}

This problem requires a substantial input from the theory of Rankin--Selberg local integrals at non-archimedean places, specifically:
\begin{enumerate}\setlength\itemsep{2pt}
\item The construction and uniqueness of essential vectors (Jacquet--Piatetski-Shapiro--Shalika, Matringe);
\item The Macdonald-type formula relating the integral with newvectors to local $L$-factors (Booker--Krishnamurthy, Kim--Krishnamurthy);
\item The Kirillov-model description of Whittaker models and the test-vector flexibility (Jacquet--Langlands, Bernstein--Zelevinsky);
\item The role of the matrix $u_Q$ as the standard ``conductor twist'' that pulls $\Pi$'s newvector into $\pi$'s conductor level.
\end{enumerate}

The author has assembled these into the proof strategy above and the assertion that the answer is YES.  A fully self-contained derivation of the test vector $V$ in Step 4 (showing the kernel $K(x,s,Q)$ is non-vanishing) is several pages and not reproduced in detail here.

\medskip
\noindent\textbf{Conjectured answer (with confidence): YES.}  The essential vector $W=W^{\rm new}$ of $\Pi$ works.  For each $\pi$, an appropriate test vector $V$ from $\pi$'s Kirillov model (depending on $c(\pi)$ and $Q$) makes the Rankin--Selberg integral a finite, nonzero entire function of $s$.

\begin{thebibliography}{9}
\bibitem{JPSS81} H.~Jacquet, I.~I.~Piatetski-Shapiro, J.~Shalika.  \emph{Conducteur des représentations du groupe linéaire}.  Math.\ Ann.\ \textbf{256} (1981), 199--214.
\bibitem{Matringe} N.~Matringe.  \emph{Essential Whittaker functions for $\mathrm{GL}(n)$}.  Doc.\ Math.\ \textbf{18} (2013), 1191--1214.
\bibitem{JacquetMacdonald} H.~Jacquet.  \emph{A correction to ``Conducteur des représentations du groupe linéaire''}.  Pacific J.\ Math.\ \textbf{260} (2012), 515--525.
\bibitem{BushnellHenniart} C.~Bushnell, G.~Henniart.  \emph{The Local Langlands Conjecture for $\mathrm{GL}(2)$}.  Springer, 2006.
\bibitem{Cogdell} J.~Cogdell.  \emph{Lectures on $L$-functions, Converse Theorems, and Functoriality for $\mathrm{GL}_n$}.  Fields Inst.\ Lectures, 2007.
\end{thebibliography}

\end{document}
